\frontmatterpage
\unchapter{Abstract}

Selecting an appropriate backend software architecture is a decisive engineering and financial choice for small and medium-sized e-commerce enterprises (SMEs), which must balance rapid feature delivery against tight infrastructure budgets and small development teams. This Bachelor's thesis presents a controlled empirical comparison between a modular monolithic architecture and a distributed microservices architecture using two production-grade reference implementations: Microsoft \texttt{eShopOnWeb} and \texttt{eShopOnContainers}.

The thesis consists of five chapters. Chapter 1 reviews the theoretical foundations of software modularity, domain-driven design, and distributed systems trade-offs. Chapter 2 analyzes the operational constraints, traffic profiles, and evolution requirements of SME e-commerce platforms. Chapter 3 details the containerized experimental testbed, workload generation protocol, and five standardized functional change scenarios. Chapter 4 reports quantitative benchmark measurements and code-change propagation results across both architectures. Chapter 5 synthesizes the empirical evidence into an actionable architectural decision framework and evaluates threats to validity.

The methodological basis combines systematic literature review, static source-code dependency and change-impact analysis across five evolution scenarios, controlled load and resource benchmarking via Docker Engine and \texttt{k6}, and structured practitioner interviews with senior backend engineers.

Empirically, the modular monolith consumes 25.5\% less idle memory (903.98 MiB vs.\ 1,213.60 MiB) and 14.7\% less memory under concurrent load (1,151.35 MiB vs.\ 1,350.20 MiB), while requiring fewer cross-boundary code changes for cross-cutting features. Conversely, microservices achieve 2.65$\times$ lower 95th-percentile tail latency (12.30 ms vs.\ 32.57 ms) and isolated service deployment, demonstrating that SMEs benefit most from starting with a modular monolith and decomposing selectively only when team or domain bottlenecks emerge.

\thesisscopeEN
